% Unidad R568: introducción, líneas francesas 70--93, pp. 7--8.
% Traducción: OpenAI Codex, GPT-6.1 Sol, esfuerzo Ultra.
Queda fuera de nuestra competencia ofrecer en esta introducción una reseña histórica, siquiera sumaria, de las nociones y los resultados expuestos.
El texto solo contendrá las referencias que se consideren especialmente útiles para su comprensión, y solo indicaremos el origen de los resultados más importantes.
Al menos en su forma, los temas tratados en nuestra obra son bastante nuevos; esto explica la escasez de referencias a los padres de la geometría algebraica del siglo~\textsc{xix} y de comienzos del siglo~\textsc{xx}, cuyos trabajos solo conocemos de oídas.
Conviene, sin embargo, decir aquí unas palabras sobre las obras que han influido más directamente en los autores y han contribuido al desarrollo del punto de vista de los esquemas.
Hay que citar, en primer lugar, el trabajo fundamental (FAC) de J.-P.~Serre, que ha servido de introducción a la geometría algebraica a más de un joven interesado en ella (entre ellos, uno de los autores de este tratado), desalentado por la aridez de los clásicos \emph{Foundations} de A.~Weil~\cite{I-18}.
Allí se demuestra por primera vez que la «topología de Zariski» de una variedad algebraica «abstracta» es plenamente adecuada para aplicarle ciertas técnicas de la topología algebraica y dar lugar, en particular, a una teoría cohomológica.
Además, la definición de variedad algebraica que allí se da es la que se presta de manera más natural a la extensión de esta noción que desarrollamos aquí\footnote{Como nos ha señalado J.-P.~Serre, conviene recordar que la idea de definir la estructura de variedad mediante un haz de anillos se debe a H.~Cartan, quien tomó esta idea como punto de partida de su teoría de los espacios analíticos.
Por supuesto, al igual que en geometría algebraica, sería importante admitir, en «geometría analítica», los elementos nilpotentes en los anillos locales de los espacios analíticos.
H.~Grauert ha abordado recientemente esta extensión de la definición de H.~Cartan y J.-P.~Serre~\cite{I-5}, y cabe esperar que pronto aparezca una exposición sistemática de geometría analítica en este marco general.
Es evidente, por otra parte, que las nociones y técnicas desarrolladas en este tratado siguen teniendo sentido en geometría analítica, aunque en esta última teoría hay que prever dificultades técnicas más considerables.
Cabe prever que la geometría algebraica, gracias a la sencillez de sus métodos, podrá servir como una especie de modelo formal para futuros desarrollos de la teoría de los espacios analíticos.}.
El propio Serre había observado, además, que la teoría cohomológica de las variedades algebraicas afines podía trasladarse sin dificultad sustituyendo las álgebras afines sobre un cuerpo por anillos conmutativos cualesquiera.
Por tanto, los capítulos~I y~II de este tratado y los dos primeros párrafos del capítulo~III pueden considerarse, en lo esencial, transposiciones fáciles, a este marco ampliado, de los principales resultados de (FAC) y de un artículo posterior del mismo autor~\cite{I-15}.
También hemos obtenido gran provecho del \emph{Séminaire de Géométrie algébrique} de C.~Chevalley~\cite{I-1}; en particular, el uso sistemático de los «conjuntos construibles» introducidos por él ha resultado muy útil en la teoría de esquemas (véase el cap.~IV).
También hemos tomado de él el estudio de los morfismos desde el punto

\oldpage[I]{8}
de vista de la dimensión (cap.~IV), que se traslada al marco de los esquemas sin cambios apreciables.
Conviene señalar, además, que la noción de «esquemas de anillos locales», introducida por Chevalley, se presta de manera natural a una extensión de la geometría algebraica (que, sin embargo, no tiene toda la flexibilidad y generalidad que pretendemos darle aquí); para la relación entre esta noción y nuestra teoría, véase el capítulo primero, \S~8.
M.~Nagata ha desarrollado una extensión de este tipo en una serie de memorias~\cite{I-9} que contienen numerosos resultados especiales sobre geometría algebraica sobre anillos de Dedekind\footnote{Entre los trabajos que se aproximan a nuestro punto de vista en geometría algebraica, señalemos el importante trabajo de E.~Kähler~\cite{I-22} y una nota reciente de Chow e Igusa~\cite{I-3}, que retoman, en el marco de la teoría de Nagata-Chevalley, ciertos resultados de (FAC) y dan también una fórmula de Künneth.}.

\bigskip
\begin{center}*\quad*\quad*\end{center}
\bigskip
